\documentclass[10pt,a4paper]{amsart}
\usepackage[margin=19mm]{geometry}
\usepackage{amsmath,amssymb,mathtools}
\usepackage[scr=rsfs,cal=euler]{mathalfa}
\usepackage{xfrac}
\usepackage[unicode,hidelinks,bookmarks=false]{hyperref}
\newcommand{\ie}{i.e.\ }
\newcommand{\cf}{cf.\ }
\newcommand{\resp}{respectively\ }
\newcommand{\Subsection}{\subsection}
\providecommand{\nobreakdash}{\nobreak\hbox{-}}
\setlength{\emergencystretch}{2em}
\multlinegap0pt
\usepackage{bm}
\usepackage{bbm}
\usepackage{dsfont}
\usepackage{xspace}
\usepackage{cancel}
\usepackage{mathtools}
\usepackage{amsthm}
\makeatletter
\@ifundefined{theorem}{\newtheorem{theorem}{Theorem}}{}
\@ifundefined{lemma}{\newtheorem{lemma}[theorem]{Lemma}}{}
\makeatother
\title[NestQuant: Nested Lattice Quantization for Matrix Products a]{NestQuant: Nested Lattice Quantization for Matrix Products and LLMs}
\author{Semyon Savkin, Eitan Porat, Or Ordentlich, Yury Polyanskiy}
\date{Source: arXiv:2502.09720v3 (CC BY 4.0)}
\begin{document}
\maketitle

\noindent\emph{Source and licence.} NestQuant: Nested Lattice Quantization for Matrix Products and LLMs. Version arXiv:2502.09720v3 (CC BY 4.0), distributed under the Creative Commons Attribution 4.0 International licence, as recorded in the arXiv licence metadata for this exact version and shown on the abstract page https://arxiv.org/abs/2502.09720v3. Full rights notice: licenses/arxiv-2502.09720-CC-BY-4.0.txt.\par\smallskip

\noindent\textit{Editorial scope.} Counted unit: the Lemma whose source label is \texttt{lemma:nestquantm} in the article (source0.tex lines 929--932, source label \texttt{lemma:nestquantm}), delivered together with the article's own complete original proof of it (source0.tex lines 934--966, one \texttt{proof} environment of 33 lines). No mathematics is written, repaired or re-derived: statement and proof are the article's own text line for line. Required context retained uncounted: none. Every definition, display and label of those lines is retained unchanged. Omitted from this package and not redistributed: 1--928, 933--933, 967--1366. Nothing inside a retained excerpt is omitted or altered: every retained body is delivered exactly as the article's lines are written, so no omission receipt of any kind accompanies this package. Citations: the retained excerpts carry no citation command at all, so no bibliography entry of the article is transcribed and no reference is redistributed. Reference adaptation: none was needed and none was made; every reference inside the delivered unit resolves inside it, so no unresolved reference or citation is produced. Printed slips kept literal: no printed word, symbol, hypothesis or number of the counted statement or of its proof was altered. Layout and class: the article's own class is \texttt{article} with packages including microtype, graphicx, booktabs, xcolor, algorithm, algpseudocode, subfigure, threeparttable, mathtools, svg, PRIMEarxiv, natbib, adjustbox, inputenc; the wrapper is the lane's pinned \texttt{amsart} editorial wrapper at 10pt on A4, which restates the theorem-like environments the retained text needs and adds the article's own macro definitions that the retained text needs (none), with extra wrapper packages (bm, bbm, dsfont, xspace, cancel, mathtools). The delivered numbering is the unit's own. Assets: no retained span contains a figure environment, a graphics-inclusion command, a PDF-page inclusion, a table or a TikZ picture; no image file is copied into this package, the package declares no asset dependency and no image file is redistributed with it. Licence: version arXiv:2502.09720v3 of this article is distributed under the Creative Commons Attribution 4.0 International licence, as recorded in the arXiv licence metadata for this exact version and shown on the abstract page https://arxiv.org/abs/2502.09720v3. A full rights notice is delivered at licenses/arxiv-2502.09720-CC-BY-4.0.txt. Characters: every retained excerpt is pure ASCII, so the delivered engine, XeLaTeX with its default Latin Modern text font, reports no missing character.
\begin{lemma}
    For a vector $x \in \mathbb{R}^8$ and a vector $v \in E_8$, $f(x + v) = f(x) + v$.
    \label{lemma:nestquantm}
\end{lemma}

\begin{proof}
    Recall that $D_8 \subset \mathbb{Z}^8$ contains all the integer vector with even sum of coordinates. The modified Gosset oracle $f$ uses modified $D_8$ oracle $g$, and at input $x$ constructs two candidate points $c_1(x) \in D_8$ and $c_2(x) \in D_8 + \frac{1}{2}$. Then, the algorithm chooses the closest point among these candidates to $x$. Note that $c_1(x) = g(x)$ and $c_2(x) = g\left(x - \frac{1}{2}\right) + \frac{1}{2}$.

    We will prove that for $u \in D_8$, $g(x + u) = g(x) + u$ for any $x \in \mathbb{R}^8$. Now, let's show the original lemma. Note that since we are choosing the closest candidate point, if the condition on candidate sets $\{c_1(x + v), c_2(x + v)\} = \{c_1(x), c_2(x)\} + v$ holds, then $f(x+v) = f(x) + v$. Now, consider two cases:

    \begin{enumerate}
        \item $v \in D_8$. Then:

        \begin{align*}
            c_1(x + v) &= g(x + v) = g(x) + v = c_1(x) + v \\
            c_2(x + v) &= g\left(x+v-\frac{1}{2}\right) + \frac{1}{2} =
            g\left(x-\frac{1}{2}\right)+\frac{1}{2}+v =
            c_2(x) + v
        \end{align*}

        \item $v \in D_8+\frac{1}{2}$. Then, we say that $v = u - \frac{1}{2} = w + \frac{1}{2}$ for $u, v \in D_8$.

        \begin{align*}
            c_1(x + v) &= g(x + v) = g\left(x + u - \frac{1}{2}\right) = g\left(x - \frac{1}{2}\right) + u
            = g\left(x - \frac{1}{2}\right) + \frac{1}{2} + v =
            c_2(x) + v \\
            c_2(x + v) &= g\left(x - \frac{1}{2} + v\right) + \frac{1}{2} = g(x + w) + \frac{1}{2} = g(x) + w + \frac{1}{2} = g(x) + v = c_1(x) + v
        \end{align*}
    \end{enumerate}

    Thus, in both cases the condition on candidate sets holds, and we get $f(x + v) = f(x) + v$.

    Now we show that if $u \in D_8$, $g(x + u) = g(x) + u$. Note that when evaluating $g(x + u)$, we will get the same rounding directions as in $g(x)$, since $u$ is an integer vector. Since $u$ has even sum of coordinates, the parity will also be the same. Then, our decision to flip the rounding of the first coordinate will also match as well. Therefore, the vector between original and rounded coordinates will be the same:

    \[
        g(x) - x = g(x + u) - x - u \Rightarrow g(x + u) = g(x) + u
    \]
\end{proof}

\end{document}
